\chapter{结论}
\label{ch:conclusion}
JHQ-GPU展示如何把结构化量化层次转化为有用的GPU搜索流水线，
以及为什么执行与精排工作量需要联合评估。四个研究问题得到具体且有边界的答案。
\section{四个研究问题的答案}
\label{sec:9-1}
\paragraph{RQ1：构建与容量。}
GPU训练、batch编码、IVF放置及packed数组形成可搜索常驻层次。
最大被测数据库为17.78百万向量、常驻20.44\,GiB，补证train+add为10.54--10.65\,s（\cref{sec:7-7}）。
Stella采样构建峰值22.58\,GiB，BGE在add入口达到26.08\,GiB。
独立重复构建和启动区间更完整地说明了已测容量边界。

\paragraph{RQ2：正确高效的执行。}
继承Cartesian编码允许实数精确拆表，完整键保留／合并及共享界论证说明分组选择
条件。两组表在22/24单元领先；相对lane精排参考，集成调度在batch 8的吞吐几何
平均为$2.727\times$，batch 1,000只有$1.015\times$
（\cref{sec:7-3,sec:7-6}）。有效并行因batch和工作分布而变；有限输出检查补充证明，
算术、并列与内存假设仍明确保留。

\paragraph{RQ3：精排预算与质量。}
无标签参考一致性无需学习预测器即可选择工作负载survivor工作量。
FULL在12/12单元满足所测相对recall损失目标，最大损失0.0006；
$S=128$一个单元则有40/64次越阈（\cref{sec:7-5}）。
参考分析解释有条件的预算搜索，不给总体GT保证。
10个FULL单元的一般完整键嵌套仍有条件，但十二个FULL有限网格选择均与枚举一致。
四个单元的直接校准加复用测量支持合适工作负载的收益，并保留困难工作负载的质量／时间权衡。

\paragraph{RQ4：竞争区间。}
0.97 recall的大batch E1插值比较中，JHQ在三个数据集第一，arXiv与IVF-RaBitQ
接近平手（\cref{sec:7-2}）；Vogue领先依赖分组路径。OpenAI E1使用TF32；同会话FP32敏感性在该目标仍领先归档
RaBitQ 7.9--8.6\%，为CAGRA fp32的$1.63$--$1.68\times$。
低recall由CAGRA int8领先，部分高目标由IVF-RaBitQ领先。
结果支持所测设备和网格内有用区间，显存另行报告；没有提供跨系统小batch或跨设备排名。

\section{系统原则与展望}
\label{sec:9-2}
最终原则是在明确质量目标下联合评估表示、执行与工作量。
表更小可能工作更多，分组更多可能重复工作，精排更少可能牺牲最重要近邻。
负结果解释数学复用和更低重建误差为何未必改善搜索前沿。
下一步应检验复用校准信息能否以更低总成本提供有效风险控制，以及显式残差放置
能否在可接受尾延迟下扩展容量（\cref{sec:8-4}）。这些是假设，来自已测局限，
不是当前系统已经提供的能力。
